
\subsubsection{Mechanistic Interpretation}

The positive correlation (r=0.72) between reliability and robustness on factual prompts is interpreted as consistent with pre-training memorization of factual knowledge. When models train on internet text corpora, they encode factual information in ways that enable both retrieving correct answers (reliability) and retrieving consistent answers across paraphrases (robustness). Strong memorization produces high reliability and high robustness; weak memorization produces low values for both. The mechanism specificity---factual prompts show r=0.72 where memorization is relevant, versus misinformation prompts showing r=0.28 where reasoning matters more---provides convergent evidence for this interpretation. Establishing definitive causality requires intervention studies manipulating memorization strength directly (e.g., ablating memory components, varying pre-training corpus composition). The observational evidence provides support consistent with the memorization mechanism but not causal proof.

The negative correlation (r=-0.25) between fairness and reliability is interpreted as consistent with RLHF fine-tuning creating an optimization trade-off. When models are trained to prioritize safety and demographic fairness, they hedge on socially sensitive questions: high fairness behavior produces factually evasive responses (low reliability), while low fairness behavior produces direct factual answers that may contain demographic bias variance (high reliability). Definitive causal attribution requires intervention experiments (e.g., comparing models with and without RLHF, varying safety constraint strength). The correlation evidence suggests the alignment tax mechanism but represents observational support rather than experimental proof.

\subsubsection{Limitations}

\textbf{Limitation 1: Underpowered Moderation Test (h-m3).} The prompt-type moderation hypothesis remains inconclusive due to pilot test with n=10 samples per stratum, whereas power analysis recommended $n\geq 85$ for 80\% power to detect r=0.3 at $\alpha =0.05$. Correlation estimates are unstable (standard error $\approx 0.35)$, producing wide 95\% CIs that include zero (factual: [-0.79, 0.38], misinformation: [-0.73, 0.50]). Fisher z-test (p=0.788) cannot distinguish whether moderation exists. The primary mechanisms (memorization, alignment tax) were robustly validated with adequate power (h-m1: n=343, h-m2: n=817). The moderation hypothesis is a refinement question---does mechanism strength vary by prompt type?---not a foundational claim. The underpowered pilot is an implementation gap (computational budget constraint), not hypothesis failure. The validated codebase and statistical pipeline can be scaled to $n\geq 100$ in future work. Claims: we cannot conclusively state whether correlation magnitude differs significantly between factual versus misinformation prompts.

\textbf{Limitation 2: Single Model Family (Llama-2 Only).} All experiments used Llama-2-7b-chat. Generalization to other architectures (GPT-4, Claude, Gemini) and model families (base models, instruction-tuned without RLHF) is unknown. Demonstrating coupling existence in one well-characterized model family is sufficient for proof-of-concept. The methodology is architecture-agnostic and can be applied to any generative LLM. Cross-architectural validation is a natural extension. Claims: observed correlation magnitudes (r=0.72, r=-0.25) are specific to Llama-2-7b under specified generation parameters. Patterns may generalize qualitatively (memorization creates positive coupling, RLHF creates negative trade-offs) but magnitudes may differ across architectures.

\textbf{Limitation 3: GPT-4-as-Judge Dependency.} Reliability scores use GPT-4-as-judge, introducing external model dependency. Assumption A1 ($\geq 90$\% agreement with human ground truth) follows standard practice but was not empirically validated in this study. GPT-4-as-judge enables automated scoring at scale (n=$817\times 3$ models = 2,451 evaluations). The large effect size (h-m1: r=0.72) provides robustness: even if GPT-4 introduces 20\% measurement noise, true correlation would remain strongly positive (r>0.5) and statistically significant. Future work can validate A1 via human annotation on $n\geq 100$ subsample. Claims: r=0.72 and r=-0.25 correlations may differ from ground-truth human-judged correlations by an unknown margin (likely $\pm 0.1-0.2)$. The qualitative pattern is robust, but precise magnitudes should be interpreted with measurement uncertainty.

\textbf{Limitation 4: Single Model Scale in h-m2/h-m3.} While the original design specified three Llama-2 scales (7B, 13B, 70B), experiments h-m2 and h-m3 tested only the 7B variant. Assumption A5 (correlations generalize across scales) remains unverified. Establishing correlation existence at one scale validates methodological feasibility. Scale generalization is an empirical question for future work. Claims: observed correlations (r=-0.25 fairness-reliability, h-m3 moderation test) are specific to Llama-2-7b and may exhibit scale-dependent effects.

\subsubsection{Broader Impact}

For trustworthiness research, this work establishes synchronized evaluation as a methodological paradigm, shifting focus from per-dimension scores to correlation structure analysis. Evaluation logs previously treated as independent datasets can be analyzed jointly to reveal training mechanism signatures.

For safety practitioners, the alignment tax quantification (r=-0.25) provides a metric for safety intervention cost-benefit analysis. Before deploying RLHF with stronger safety constraints, practitioners can measure baseline fairness-reliability correlation, estimate expected reliability degradation from fairness improvement, and decide whether safety gain justifies accuracy cost.

For model development, the memorization mechanism signature (r=0.72 on factual content) enables diagnostic uses: low r\_reliability-robustness on factual prompts indicates weak memorization; targeted pre-training on domain knowledge should improve both dimensions together.
